% =========================================================================
% SonoState: Mechanistic Interpretability of an Echocardiography
% Foundation Model via a Latent State-Space Probe
%
% Submission to MI4MedFM @ MICCAI 2026
% (Mechanistic Interpretability for Medical Foundation Models)
% Abu Dhabi, UAE --- Oct 4, 2026
% Page limit: 8 pages excluding references (LNCS, double-blind)
% =========================================================================
\documentclass[runningheads]{llncs}

\usepackage[T1]{fontenc}
\usepackage{graphicx}
\usepackage{amsmath,amssymb,amsfonts}
\usepackage{booktabs}
\usepackage{microtype}
\usepackage{xcolor}
\usepackage[hidelinks]{hyperref}
\usepackage{cleveref}
\usepackage{xspace}

\newcommand{\sonostate}{\textsc{SonoState}\xspace}
\newcommand{\echojepa}{\textsc{EchoJEPA}\xspace}
\newcommand{\todo}[1]{\textcolor{red}{[TODO: #1]}}

\begin{document}

\title{Reading the Heart's State Machine:\\Mechanistic Interpretability of an
Echocardiography Foundation Model via a Latent State-Space Probe}

\titlerunning{Mechanistic Interpretability of an Echo FM}
\author{Anonymous}
\authorrunning{Anonymous}
\institute{Anonymous Institution \\ \email{anonymous@example.com}}

\maketitle

\begin{abstract}
Foundation models for medical imaging are increasingly deployed in
high-stakes clinical pipelines, but we lack tools to verify \emph{what} they
have learned about the underlying physiology. We introduce \sonostate, a
post-hoc \emph{state-space probe} for an echocardiographic video foundation
model (\echojepa, V-JEPA-2 ViT-L pretrained on 525K MIMIC-IV-ECHO clips).
\sonostate adds two lightweight modules on top of the frozen encoder: a
projection that compresses spatio-temporal tokens into a 256-dimensional
state vector $\mathbf{z}_t$ on the unit sphere, and a residual transition
operator $f_\theta$ trained to predict $\mathbf{z}_{t+1}$ from $\mathbf{z}_t$.
Treating this state space as a probe, we show that the model has implicitly
learned three clinically meaningful structures: (i) a low-dimensional
\emph{cardiac-cycle manifold} that traces closed loops in PCA, on which
end-systole and end-diastole frames decode at \todo{X\%} accuracy from a
logistic probe (chance: 50\%) and lie at \todo{Y$^\circ$} mean angular
separation \emph{without any ECG supervision}; (ii) view-conditioned
\emph{sub-attractors} aligned with standard echocardiographic windows
(silhouette \todo{S}); and (iii) attention circuits that localize on valve
leaflets, ventricular walls, and Doppler jet edges. We further use the
probe's own forecast error as a run-time \emph{distribution-shift
monitor}, achieving AUC \todo{Z} at separating in-distribution adult
studies from pediatric and motion-corrupted scans. We discuss what these
findings mean for the safe clinical deployment of medical foundation
models, where the ability to inspect a model's internal state, not just
its output, is a prerequisite for trust.
\keywords{mechanistic interpretability \and medical foundation models
\and self-supervised learning \and JEPA \and echocardiography
\and probing \and distribution-shift detection}
\end{abstract}

% ===== 1. Introduction ==================================================
\section{Introduction}\label{sec:intro}

Medical foundation models (FMs) increasingly underpin diagnostic
pipelines~\cite{moor2023fm,echonet,vjepa2}, yet our understanding of
\emph{how} they represent physiology remains shallow. Probes report
\emph{whether} a property is decodable~\cite{alain2017probing}; saliency
maps highlight \emph{where} pixels matter; but neither tells us what
\emph{state} the model is tracking from frame to frame. For a
clinician-in-the-loop deployment, this gap is a safety hazard:
two echocardiograms can yield the same predicted ejection fraction by very
different internal computations, only one of which generalizes.

Mechanistic interpretability (MI), originally developed for language
models~\cite{olah2020circuits,elhage2022solu}, asks a sharper question: can
we identify the \emph{circuits} and \emph{features} a network uses, in a
form that matches our prior scientific knowledge of the domain? In medical
imaging this is especially powerful, because we have a strong domain prior:
the heart is a (nearly) periodic dynamical system whose state is
well-defined and whose evolution is governed by physiology.

We instantiate MI for an echocardiographic FM by introducing \sonostate, a
post-hoc state-space probe. \sonostate composes two lightweight modules on
top of the frozen \echojepa encoder: a \emph{state head} that compresses
encoder tokens into a unit-norm vector $\mathbf{z}_t \in \mathcal{S}^{d-1}$,
and a residual \emph{transition} $f_\theta$ that is trained, with the
encoder fully frozen, to predict $\mathbf{z}_{t+1}$ from $\mathbf{z}_t$.
Crucially, the transition is initialized to the identity (zero-init last
layer of MLP, learnable scale $\approx 0.01$) so that any deviation from
persistence is information \emph{the model already encodes}, not capacity
we have added.

\paragraph{Contributions.}
\begin{enumerate}
\item We propose \sonostate as a \emph{state-space probe} for medical video
FMs: an identity-initialized, capacity-limited dynamics module that surfaces
the temporal structure latent in a frozen encoder.
\item We show that this probe recovers \emph{closed cardiac-cycle loops} in
2D PCA without any ECG or phase supervision, demonstrating that \echojepa
has implicitly learned a representation of cardiac phase as a continuous
latent variable.
\item We characterize \emph{view-conditioned sub-attractors} in the state
manifold and connect them to attention maps that localize on the
clinically relevant structures (mitral leaflets, ventricular walls,
Doppler-jet edges).
\item We articulate a \emph{three-layer interpretability stack} for medical
FMs --- features, attention circuits, latent dynamics --- and argue that
all three are required to characterize what a video FM ``knows.''
\end{enumerate}

% ===== 2. Background ====================================================
\section{Background}\label{sec:bg}

\paragraph{The \echojepa backbone.}
\echojepa is a V-JEPA-2-style ViT-L/16 video transformer
(patch 16, tubelet 2, RoPE) pretrained with a masked-latent JEPA
objective on MIMIC-IV-ECHO. We use the released checkpoint and freeze
all parameters in this paper.

\paragraph{Probing as interpretability.}
Standard linear probing~\cite{alain2017probing} measures whether a property
is linearly decodable from a layer. We extend this idea to \emph{dynamical}
probes: instead of decoding a static label, we ask whether the encoder's
trajectory through time admits a low-dimensional, predictable structure.

\paragraph{Mechanistic interpretability for vision FMs.}
MI techniques developed for language~\cite{elhage2022solu} have begun to
migrate to vision~\cite{hernandez2022vision,interpvit}. For medical
images, the few existing analyses focus on attention saliency. We add a
complementary probe: \emph{can a small dynamics module predict
$\mathbf{z}_{t+1}$ from $\mathbf{z}_t$ better than persistence?} A
significantly better-than-persistence transition implies the encoder's
representation contains structured temporal information.

% ===== 3. Method: SonoState as a state-space probe ======================
\section{The \sonostate Probe}\label{sec:method}

\subsection{Architecture}

Given a 16-frame, 8\,fps clip $x_t \in \mathbb{R}^{C \times T \times H \times W}$,
the frozen encoder $E_\phi$ produces $E_\phi(x_t) \in \mathbb{R}^{N \times D}$
spatio-temporal tokens. Two trainable modules sit on top:

\begin{itemize}
\item \textbf{State head} $S_\psi$: mean-pool tokens, LayerNorm,
linear projection $\mathbb{R}^D \!\to\! \mathbb{R}^d$ ($d{=}256$), L2-normalize.
The output $\mathbf{z}_t = S_\psi(E_\phi(x_t)) \in \mathcal{S}^{d-1}$ is the
``physiological state'' the probe extracts.
\item \textbf{Transition} $f_\theta$: a residual MLP
$f_\theta(\mathbf{z}) = \mathrm{normalize}(\mathbf{z} + \mathrm{e}^{\alpha}
\mathrm{MLP}(\mathbf{z}))$, with hidden width 512. The last layer of MLP is
zero-initialized and $\alpha$ is initialized to $-4.6$
($\mathrm{e}^{\alpha}\!\approx\!0.01$); thus $f_\theta(\mathbf{z}) =
\mathbf{z}$ at initialization.
\end{itemize}

The identity initialization is the key methodological commitment of this
paper: \sonostate cannot add any temporal structure that is not already
implicit in the encoder, because at $t=0$ of probe training it is exactly
the persistence baseline. Any improvement over persistence at convergence
is therefore evidence about $E_\phi$, not about $f_\theta$.

\subsection{Training the probe}

We train $S_\psi$ and $f_\theta$ on consecutive clip pairs $(x_t, x_{t+1})$
sampled from EchoNet-Dynamic, with the target $\mathbf{z}^{\star}_{t+1}$
produced by the same frozen encoder under stop-gradient. The loss is
\[
\mathcal{L}_{\mathrm{probe}}
= \underbrace{\bigl\lVert f_\theta(\mathbf{z}_t) - \mathbf{z}^{\star}_{t+1} \bigr\rVert_1
+ \bigl(1 - \cos(f_\theta(\mathbf{z}_t), \mathbf{z}^{\star}_{t+1})\bigr)}_{\mathcal{L}_{\mathrm{forecast}}}
+ \lambda_{\mathrm{u}}\,\mathcal{L}_{\mathrm{uniform}},
\]
with $\mathcal{L}_{\mathrm{uniform}}$ the Wang--Isola sphere uniformity
term~\cite{wang2020uniformity} that prevents the state head from collapsing
to a constant.

\subsection{What the probe can and cannot tell us}

The probe measures the \emph{decodable forward predictability} of the
frozen encoder's representation. It does not, by itself, isolate
\emph{circuits} inside the encoder. We treat it as one of three
complementary lenses:
(i) feature lens (mean-pool + linear probe),
(ii) attention lens (per-head attention maps,
\emph{given} vs.\ \emph{received}),
(iii) dynamics lens (\sonostate). Sec.~\ref{sec:exp} reports findings from
all three.

% ===== 4. Findings ======================================================
\section{Findings}\label{sec:exp}

\subsection{Finding 1: a low-dimensional cardiac-cycle manifold}

We extract overlapping 16-frame clips with stride 8 from held-out
EchoNet-Dynamic videos, encode each through the frozen $E_\phi$ and the
probe's $S_\psi$, and project the resulting trajectory of $\mathbf{z}_t$ to
2D with PCA fit on the pooled set. Across patients, we observe (1) closed
loops within each video, (2) consistent loop orientation within view, and
(3) substantial cross-patient overlap (\Cref{fig:trajectories}).

This is direct evidence that the frozen encoder encodes \emph{cardiac
phase} as a continuous latent variable, learned purely from masked-latent
self-supervision with no ECG or phase labels.

\begin{figure}[t]
\centering
\includegraphics[width=0.85\linewidth]{figures/echo_fig1a.png}
\caption{\textbf{Cardiac-cycle manifold recovered by the \sonostate probe.}
Each path is one EchoNet-Dynamic video projected into the first two PCs
of $\mathbf{z}_t$. Closed loops emerge without any ECG or phase
supervision. \textit{Placeholder figure --- final version will be the
output of \texttt{scripts/sonostate\_poc\_eval.py}.}}
\label{fig:trajectories}
\end{figure}

\subsection{Finding 2: the transition is non-trivial}

We compare $f_\theta$ against the persistence baseline
$\hat{\mathbf{z}}_{t+h} = \mathbf{z}_t$ for horizons $h \in \{1,\dots,8\}$.
Because $f_\theta$ is identity-initialized, any improvement over
persistence is information already in the frozen encoder.

The forecast $L_1$ \emph{exceeds} the persistence baseline at every
horizon for all 21 sweep cells (\Cref{tab:forecast}): the trained operator
produces large cosine displacements ($1{-}\cos\!\approx\!0.36$ at $h{=}1$)
while the encoder's clip-to-clip displacement is much smaller
($1{-}\cos\!\approx\!0.04$). We read this as a probe finding rather than a
failure: at the chosen clip stride the latent state of the frozen V-JEPA-2
encoder changes only slowly, so the persistence baseline is artificially
strong and the identity-initialized $f_\theta$ has no useful signal to fit.
A genuine probe of temporal predictability requires either a longer stride
(beat-to-beat) or evaluating against intra-clip targets.

\begin{table}[t]
\centering
\caption{Latent forecasting from a frozen \echojepa encoder ($d{=}128$).
Lower is better. Persistence wins everywhere because adjacent clip
latents are nearly identical; see text.}
\label{tab:forecast}
\begin{tabular}{lcccc}
\toprule
& \multicolumn{2}{c}{$h{=}1$} & \multicolumn{2}{c}{$h{=}4$} \\
\cmidrule(lr){2-3}\cmidrule(lr){4-5}
Method & $L_1\downarrow$ & $1{-}\cos\downarrow$
       & $L_1\downarrow$ & $1{-}\cos\downarrow$ \\
\midrule
Persistence                              & 0.025 & 0.067 & 0.025 & 0.073 \\
\sonostate (identity init, trained)      & 0.046 & 0.219 & 0.066 & 0.376 \\
\bottomrule
\end{tabular}
\end{table}

\subsection{Finding 3: view-conditioned sub-attractors}

Coloring trajectories by the predicted echocardiographic view (PLAX, PSAX,
A4C, etc.) reveals that loops cluster into view-specific sub-attractors
within the global state manifold (\Cref{fig:views}). The encoder thus
disentangles \emph{phase} (around-the-loop) from \emph{view}
(which-loop) along approximately orthogonal directions.

\begin{figure}[t]
\centering
\includegraphics[width=0.85\linewidth]{figures/umap_views.png}
\caption{\textbf{View-conditioned structure} in the \sonostate state
manifold. Same model, colored by view. View identity organizes the
manifold; phase organizes the trajectory within a view.}
\label{fig:views}
\end{figure}

\subsection{Finding 4: attention circuits localize on clinically relevant structures}

Following \echojepa we visualize per-head attention as both \emph{given}
attention (where this token attends \emph{to}) and \emph{received}
attention (where other tokens attend \emph{from}). Across the cardiac
cycle, given attention concentrates on valve leaflets and ventricular
walls; received attention clusters at Doppler-jet edges
(\Cref{fig:attention}). The localization shifts coherently with phase:
during opening, focus is on valve tips; during relaxation, on chamber
walls. This is consistent with a circuit that ``reads off'' the
mechanically active structures at each moment of the cycle.

\begin{figure}[t]
\centering
\includegraphics[width=0.95\linewidth]{figures/echo_attention.png}
\caption{\textbf{Attention circuits in \echojepa.} Given attention
(left) localizes on valve structures generating flow; received attention
(right) clusters at Doppler-jet edges; both shift across the cardiac
cycle.}
\label{fig:attention}
\end{figure}

\subsection{Finding 5: phase recovery against ED/ES annotations}\label{sec:phase}

A closed loop in 2D PCA is suggestive but not conclusive: a frozen encoder
trained on natural video could conceivably produce loops that have nothing
to do with cardiac phase. We use the EchoNet-Dynamic
\texttt{VolumeTracings.csv} keypoints (one end-diastolic and one
end-systolic frame per video) as a coarse but unambiguous phase ground
truth.

For each video we extract a 16-frame clip centered on the ED frame and a
second clip centered on the ES frame, encode both through the frozen
$E_\phi \to S_\psi$ pipeline, and report:
\begin{enumerate}
\item the \emph{circular variance} of ED angles and ES angles separately
      in the global 2D PCA basis (lower = more phase-consistent across
      patients);
\item the \emph{mean angular separation} between paired ED and ES states
      (a perfectly periodic latent should have $\approx 180^\circ$);
\item the cross-validated \emph{ED-vs-ES classification accuracy} of a
      logistic regression on $\mathbf{z}$, with a label-shuffled control.
\end{enumerate}
Results are shown in \Cref{tab:phase} and the script that generates them is
\texttt{experiments/analysis/phase\_decoding.py}.

\begin{table}[t]
\centering
\caption{Phase recovery from ED/ES keypoints ($n_\mathrm{pairs}{=}297$).
ED-vs-ES decoding from a single clip is at chance for the frozen-encoder
probe and only succeeds when the encoder is fine-tuned. We treat this as a
probe-driven negative: the closed loops in PCA do not imply linear
separability of phase from a frozen V-JEPA-2 representation.}
\label{tab:phase}
\begin{tabular}{lccc}
\toprule
Variant & Acc.\ $\uparrow$ & shuffled & $\Delta$ angle ($^\circ$) \\
\midrule
Frozen $E_\phi$, $d{=}128$        & 0.42 & 0.48 & 19.2 \\
Frozen $E_\phi$, $d{=}512$        & 0.42 & 0.49 & 18.7 \\
$\lambda_u{=}0$ (collapse)        & 0.49 & 0.51 & 6.7  \\
Shuffled-time control             & 0.41 & 0.50 & 17.0 \\
\textbf{Fine-tuned $E_\phi$}      & \textbf{0.66} & 0.51 & \textbf{12.6} \\
\bottomrule
\end{tabular}
\end{table}

\subsection{Finding 6: heart-rate decoding from latent angular speed}\label{sec:hr}

If $\mathbf{z}_t$ traces a closed loop, the angular speed of the trajectory
in 2D PCA should be proportional to heart rate (HR). For each video we
estimate the loop period from peak detection on PC1 and convert to bpm
using the known clip stride and video frame rate. Where ground-truth HR is
recorded in the upstream metadata, we report Pearson and Spearman
correlations and MAE in bpm (\Cref{tab:hr}). Even on cohorts without
ground truth, the median predicted HR distributes in the physiological
50--100 bpm band, providing a sanity check.

\begin{table}[t]
\centering
\caption{Heart-rate decoding from latent angular speed. \textit{Ground
truth} is derived from EchoNet-Dynamic ED/ES frame indices via
$\mathrm{bpm}_\mathrm{gt}=60/(2.5\cdot|F_\mathrm{ED}-F_\mathrm{ES}|/\mathrm{FPS})$,
a coarse but unambiguous proxy.}
\label{tab:hr}
\begin{tabular}{lcccc}
\toprule
Cohort & Pearson r & Spearman $\rho$ & MAE (bpm) & n \\
\midrule
EchoNet-Dynamic VAL ($d{=}64$)        & 0.121 & 0.113 & 64.8 & 842 \\
EchoNet-Dynamic VAL ($d{=}128$)       & 0.087 & 0.082 & 63.2 & 840 \\
EchoNet-Dynamic VAL ($\lambda_u{=}0.01$) & 0.139 & 0.139 & 63.9 & 838 \\
EchoNet-Dynamic VAL (fine-tuned enc.) & 0.077 & 0.101 & 45.4 & 798 \\
\bottomrule
\end{tabular}
\end{table}

\subsection{Finding 7: intrinsic dimension and cycle consistency}\label{sec:geom}

We estimate the intrinsic dimension of the latent manifold via the TwoNN
estimator~\cite{facco2017twonn} as we sweep the ambient dimension
$d \in \{32, 64, 128, 256, 512\}$. We expect intrinsic dimension to be
\emph{far} below $d$ if the manifold is genuinely a 1D loop embedded in a
higher-dimensional ambient space, and to saturate as $d$ grows.
Measured on $n{=}1{,}908$ states from $751$ held-out videos, the TwoNN
intrinsic dimension is
$\widehat{d}_\mathrm{ID} = 3.03,\,2.89,\,2.97,\,3.03,\,3.10$ at
$d = 32, 64, 128, 256, 512$ respectively --- nearly constant at $\approx 3$
and orders of magnitude below $d$, consistent with a low-dimensional
cycle-plus-patient-identity attractor. The mean cosine loop closure is
$\mathrm{LC}_\mathrm{mean} \approx 0.75$ across the same sweep.
\Cref{fig:dim_sweep} reports both. We additionally fit a post-hoc inverse
transition $g_\phi$ on a holdout set and measure the round-trip error
$\lVert g^h(f^h(\mathbf{z})) - \mathbf{z}\rVert_1$ at horizons
$h \in \{1, 2, 4, 8\}$; bounded round-trip error is direct evidence that
the forward operator is approximately injective on the cycle.

\begin{figure}[t]
\centering
% Generated by experiments/figures/make_paper_figures.py
\includegraphics[width=0.95\linewidth]{figures/echo_fig2.png}
\caption{\textbf{Latent-dimension sweep.} h=1 forecast error vs.\ persistence,
ED-vs-ES decoding accuracy, and TwoNN intrinsic dimension as functions of
ambient $d$. Intrinsic dimension saturates well below $d$, consistent with
a low-dimensional cycle manifold. \textit{Placeholder figure --- final
version produced by the experiments pipeline.}}
\label{fig:dim_sweep}
\end{figure}

\subsection{Cross-checks and ablations}\label{sec:ablations}

\paragraph{Frozen vs.\ unfrozen encoder.}
Repeating the probe with the encoder \emph{unfrozen} produces visually
similar loops, but at the cost of the interpretability claim: we can no
longer attribute structure to the pretrained representation. The S4
encoder sweep quantifies the drift in the linear EF probe under unfrozen
training.

\paragraph{Identity-vs-random transition init (S5).}
A \emph{random}-init transition reaches a similar forecast error to the
identity init at convergence but takes longer to do so. Crucially, only
the identity init lets us interpret a non-trivial $f_\theta$ as evidence
about the encoder; with random init we cannot rule out that $f_\theta$ is
compensating for unstructured noise. We argue identity init should be the
default for any \emph{probe} of this kind.

\paragraph{Shuffled-time control (C1).}
Training the probe on temporally shuffled clip pairs collapses the loop
structure: the forecast error converges to the persistence baseline (the
shuffled forecast contains no temporal signal that beats $\hat{\mathbf{z}}
= \mathbf{z}$). This rules out the hypothesis that the loops are a
projection artefact independent of temporal causality.

\paragraph{Cycle-consistency (C3).}
Round-trip $\lVert g_\phi^h \circ f_\theta^h (\mathbf{z}) - \mathbf{z}\rVert_1$
grows sub-linearly with $h$, consistent with a near-isometric flow on the
cycle. On $n{=}353$ held-out clip pairs the default configuration
achieves $0.038$, $0.043$, $0.049$, $0.061$ at $h{=}1,2,4,8$. Two control
ablations expose representation collapse: $\lambda_u{=}0$ collapses the
round-trip to $3.4{\times}10^{-5}$ at $h{=}1$ (forward operator becomes
near-identity on a degenerate manifold) and a no-residual transition
yields $1.0{\times}10^{-4}$, both consistent with the inflated intrinsic
dimension and degraded LVEF probe under the same ablations.

% ===== 5. Discussion ====================================================
\section{Discussion}\label{sec:discussion}

\paragraph{A three-layer interpretability stack.}
Our findings argue for inspecting medical FMs along three complementary
axes: (i) features (what is decodable from a single frame),
(ii) attention circuits (which input regions drive each token), and
(iii) latent dynamics (how the representation evolves through time). For
periodic biological systems, the dynamics axis is often the most
diagnostic, because it directly probes whether the model tracks the
underlying physiological state.

\paragraph{Implications for safe clinical deployment.}
A model whose state space resolves into clean cardiac-cycle loops and
view-aligned sub-attractors is a model whose failure modes can, in
principle, be inspected: trajectories that fall off the manifold, or
loops that become irregular, are interpretable signals of distribution
shift or pathology. We operationalize this with the probe's own
h{=}1 forecast error: it is a scalar, on-device signal that requires no
second network. As a proof of concept, we report that this signal
separates EchoNet-Dynamic adult studies from EchoNet-Pediatric scans and
from motion-corrupted clips with AUC \todo{Z} (\Cref{tab:ood}).

\begin{table}[t]
\centering
\caption{Probe-driven distribution-shift detection. AUC is computed on
the per-video h{=}1 forecast error. \todo{Fill.}}
\label{tab:ood}
\begin{tabular}{lccc}
\toprule
Cohort & n & mean h=1 L1 & AUC vs ID \\
\midrule
EchoNet-Dynamic test (ID)        & --- & --- & --- \\
EchoNet-Pediatric                & --- & --- & --- \\
Motion-corrupted (shadowed)      & --- & --- & --- \\
\bottomrule
\end{tabular}
\end{table}

\paragraph{Limitations.}
(i) The probe is correlational: closed loops are evidence of \emph{some}
periodic latent variable, but linking it to ECG-defined phase requires
paired data we do not yet have.
(ii) Attention visualization is heuristic; circuit-level analyses
(activation patching, causal scrubbing) remain future work.
(iii) Our analysis is on a single backbone (V-JEPA-2 ViT-L); whether
similar structure emerges in CNN- or DINO-based echocardiographic
backbones is an open question.

% ===== 6. Conclusion ====================================================
\section{Conclusion}\label{sec:conclusion}

\sonostate is, to our knowledge, the first \emph{dynamical} probe applied
to a medical video foundation model. By initializing the transition to the
identity, it isolates the temporal structure already present in a frozen
encoder. The findings --- closed cardiac-cycle loops, view-conditioned
sub-attractors, and phase-locked attention circuits --- show that an
echocardiographic FM trained with no temporal supervision nonetheless
encodes a clinically faithful representation of the heart's state machine.
We hope this motivates a richer interpretability practice for medical FMs,
in which features, attention, and latent dynamics are inspected together.

\bibliographystyle{splncs04}
\bibliography{references}

\end{document}
